% Generated by scripts/assemble.py; corrections belong in source/corrections.json.
\examhead{2}{北京师范大学 2025 年数学分析试卷}{9}
% p010-b18; source PDF page 10
\blockanchor{p010-b18}
\question{1}(15 分) 已知数列 $\{a_{n}\}$ 满足 $a_{1} = \sqrt{3}, a_{n+1} = \sqrt{3 + a_{n}}$, 证明: $\{a_{n}\}$ 收敛.\par
% p010-b19; source PDF page 10
\blockanchor{p010-b19}
\question{2}(15 分) 求极限 $\lim_{n\to \infty}\frac{\pi}{n}\sum_{k = 1}^{n}\frac{1}{2 + \cos\frac{k\pi}{n}}.$\par
% p010-b20; source PDF page 10
\par\Needspace{6\baselineskip}
\blockanchor{p010-b20}
\question{3}(15 分) 已知 $f(x)$ 在 $[0,2]$ 上可导, 且 $f(0) = 0$, $f'(x)$ 在 $[0,2]$ 上连续. 证明:\par
\nopagebreak[4]
% p010-b21; source PDF page 10
\blockanchor{p010-b21}
\displayformula{\int_ {0} ^ {2} | f (x) f ^ {\prime} (x) | \mathrm{d} x \leq \int_ {0} ^ {2} | f ^ {\prime} (x) | ^ {2} \mathrm{d} x.}
% p010-b22; source PDF page 10
\blockanchor{p010-b22}
\question{4}(15 分) 求级数和 $\sum_{n=0}^{\infty} \frac{1}{n!}(n+1)(n+2)x^{n}$.\par
% p011-b00; source PDF page 11
\par\Needspace{6\baselineskip}
\blockanchor{p011-b00}
\question{5}(15 分) 已知 $A \in \mathbb{R}$, 函数 $f(x)$ 在 $[0, +\infty)$ 上连续, 且 $\lim_{x \to +\infty} f(x) = A$. 证明:\par
\nopagebreak[4]
% p011-b01; source PDF page 11
\blockanchor{p011-b01}
\displayformula{\lim _ {x \rightarrow + \infty} \frac {1}{x} \int_ {0} ^ {x} f (t) \mathrm{d} t = A.}
% p011-b02; source PDF page 11
\blockanchor{p011-b02}
\question{6}(15 分) 判断 $\int_{2}^{+\infty} \frac{\sin x}{\ln x} \arctan x \, \mathrm{d}x$ 的敛散性, 并说明理由.\par
% p011-b03; source PDF page 11
\par\Needspace{6\baselineskip}
\blockanchor{p011-b03}
\question{7}(15 分) 求\par
\nopagebreak[4]
% p011-b04; source PDF page 11
\blockanchor{p011-b04}
\displayformula{f (x, y) = - x ^ {2} y (x + y - 4)}
% p011-b05; source PDF page 11
\blockanchor{p011-b05}
在与 $x$ 轴, $y$ 轴, 直线 $y = 6 - x$ 围成的闭区域 $D$ 中的极值, 最大值与最小值.\par
% p011-b06; source PDF page 11
\blockanchor{p011-b06}
\question{8}(15 分) 已知 $f(x, y) = g(|xy|)$，且 $g(0) = 0$，在 $z > 0$ 附近，有 $|g(z)| \leq z^{\alpha}$，$\alpha > \frac{1}{2}$. 证明：$f(x, y)$ 在 (0,0) 处可微.\par
% p011-b07; source PDF page 11
\par\Needspace{6\baselineskip}
\blockanchor{p011-b07}
\question{9}(15 分) 已知 $\Sigma$ 为单位球面 $x^{2} + y^{2} + z^{2} = 1$, 取外侧. 求曲面积分\par
\nopagebreak[4]
% p011-b08; source PDF page 11
\blockanchor{p011-b08}
\displayformula{\iint_ {\Sigma} \frac {x \mathrm{d} y \mathrm{d} z + y \mathrm{d} z \mathrm{d} x + z \mathrm{d} x \mathrm{d} y}{(2 x ^ {2} + 2 y ^ {2} + z ^ {2}) ^ {\frac {3}{2}}}.}
% p011-b09; source PDF page 11
\par\Needspace{6\baselineskip}
\blockanchor{p011-b09}
\question{10}(15 分) 求三重积分\par
\nopagebreak[4]
% p011-b10; source PDF page 11
\blockanchor{p011-b10}
\displayformula{\iiint_ {x ^ {2} + y ^ {2} + z ^ {2} \leq 2 z} \left(\sqrt {x ^ {2} + y ^ {2} + z ^ {2}} + x - y ^ {3}\right) \mathrm{d} V.}
